\begin{ouverture}[courbe]{7}{Proportionnalité et pourcentages}{Partie 2 · Algèbre et modèles}
{Passage à l’unité · Quatrième proportionnelle · Échelles}
\objectif{Reconnaître une relation proportionnelle.}
\objectif{Calculer une quatrième proportionnelle avec plusieurs méthodes.}
\objectif{Utiliser pourcentages, échelles et formules.}
\prerequis{Fractions et calculs décimaux.}
\end{ouverture}
\pagestyle{otMain}\essentiel\fontsize{11}{15}\selectfont
\phantomsection\label{cours-7-1}
\essSection{1}{Reconnaître et calculer}
\begin{definition}Deux grandeurs sont proportionnelles si l’on passe des valeurs de la première à celles de la seconde en multipliant toujours par le même nombre, le coefficient de proportionnalité. Les unités et le sens du passage doivent être précisés.\end{definition}
\begin{exemple}Si 4 cahiers identiques coûtent 10 euros sans frais fixes, un cahier coûte $10/4=2{,}50$ euros. Pour 7 cahiers, on paie $7\times2{,}50=17{,}50$ euros. On peut aussi calculer directement $7\times10/4$.\end{exemple}
\begin{propriete}Dans un tableau proportionnel, on peut multiplier une colonne par un nombre ou additionner des colonnes. Pour deux couples $(a,b)$ et $(c,d)$ proportionnels, $ad=bc$ : les produits en croix sont égaux.\end{propriete}
\begin{attention}Une hausse régulière ne suffit pas : un tarif de 5 euros fixes plus 2 euros par article n’est pas proportionnel. Des points alignés doivent aussi être alignés avec l’origine pour représenter une proportionnalité.\end{attention}
\clearpage
\phantomsection\label{cours-7-2}
\essSection{2}{Appliquer dans plusieurs contextes}
\begin{propriete}Prendre $p\,\%$ d’une quantité revient à la multiplier par $p/100$. Pour trouver un pourcentage, diviser la partie par le total, puis multiplier par 100. Une baisse de 20\,\% retire 20\,\% de la valeur de départ de cette baisse.\end{propriete}
\begin{definition}Une échelle est le rapport longueur représentée sur longueur réelle, mesurées dans la même unité. À l’échelle $1/200$, 1 cm sur le plan représente 200 cm en réalité.\end{definition}
\begin{exemple}Sur 80 élèves, 28 prennent le bus : $28/80=0{,}35$, soit 35\,\%. Un article de 60 euros baissé de 15\,\% perd $0{,}15\times60=9$ euros et coûte donc 51 euros.\end{exemple}
\begin{savoirfaire}[Lire une formule]Si une résistance vaut $R=8\ \Omega$ et si $U=RI$, alors $U=8I$ : tension et intensité sont proportionnelles dans ce modèle. Pour un cercle, $L=\pi d$ relie aussi deux grandeurs proportionnelles, la circonférence et le diamètre.\end{savoirfaire}
\begin{attention}Pour deux pourcentages successifs, identifier la nouvelle base à chaque étape. On ne les additionne pas automatiquement. Les méthodes de ce chapitre reposent sur des parts et des tableaux ; le vocabulaire de fonction linéaire n’est pas nécessaire.\end{attention}

\clearpage\rubrique{Découverte}
\begin{blocLivre}{ex-7-1}\begin{exercice}[1]{Passer à l’unité}
Six bouteilles identiques coûtent 9 euros. Combien coûtent dix bouteilles, sans frais fixes ?
\end{exercice}\end{blocLivre}
\begin{blocLivre}{ex-7-2}\begin{exercice}[1]{Compléter un tableau}
Le prix est proportionnel à la masse. Compléter le tableau.\begin{center}\renewcommand{\arraystretch}{1.35}\begin{tabular}{c|ccc}\toprule
Masse (kg) & 2 & 5 & 8\\\midrule
Prix (euros) & 7 & ? & ?\\
\bottomrule\end{tabular}\end{center}
\end{exercice}\end{blocLivre}
\begin{blocLivre}{ex-7-3}\begin{exercice}[1]{Une part en pourcentage}
Calculer 18\,\% de 250.
\end{exercice}\end{blocLivre}
\begin{blocLivre}{ex-7-4}\begin{exercice}[1]{Retrouver le pourcentage}
Dans un groupe de 32 élèves, 12 sont demi-pensionnaires. Quel pourcentage cela représente-t-il ?
\end{exercice}\end{blocLivre}
\begin{blocLivre}{ex-7-5}\begin{exercice}[1]{Lire une échelle}
Sur un plan à l’échelle $1/500$, un mur mesure 3,6 cm. Quelle longueur réelle représente-t-il en mètres ?
\end{exercice}\end{blocLivre}

\rubrique{Entraînement}
\begin{blocLivre}{ex-7-6}\begin{exercice}[2]{Une quatrième proportionnelle}
Huit pièces identiques pèsent 14,4 kg. Quelle est la masse de cinq pièces ? Donner deux méthodes.
\end{exercice}\end{blocLivre}
\begin{blocLivre}{ex-7-7}\begin{exercice}[2]{Une remise}
Un sac coûte 80 euros. On lui applique une remise de 12\,\%. Quel est son nouveau prix ?
\end{exercice}\end{blocLivre}
\begin{blocLivre}{ex-7-8}\begin{exercice}[2]{Un tarif non proportionnel}
Un service facture 4 euros fixes puis 3 euros par heure. Calculer pour 1 h et 2 h. Le prix est-il proportionnel au temps ?
\end{exercice}\end{blocLivre}
\begin{blocLivre}{ex-7-9}\begin{exercice}[2]{Un graphique}
Le graphique représente un prix en fonction d’une masse. Reconnaître la proportionnalité et lire le prix de 3 kg.
\label{fig-ex-7-9}\figureGeo{proportion-prix}
\end{exercice}\end{blocLivre}
\begin{blocLivre}{ex-7-10}\begin{exercice}[2]{Une formule en sciences}
Dans le modèle $U=RI$, $R=15\ \Omega$. Calculer $U$ pour $I=0{,}4$ A, puis $I$ quand $U=9$ V.
\end{exercice}\end{blocLivre}

\rubrique{Problèmes}
\begin{blocLivre}{ex-7-11}\begin{exercice}[3]{Deux remises successives}
Un article de 100 euros subit une baisse de 20\,\%, puis de 10\,\% du nouveau prix. Quel prix obtient-on ? La baisse totale vaut-elle 30\,\% ?
\end{exercice}\end{blocLivre}
\begin{blocLivre}{ex-7-12}\begin{exercice}[3]{Une recette et des lots}
Une recette pour 6 personnes demande 450 g de farine. On la prépare pour 14 personnes. La farine est vendue en paquets de 500 g. Combien de paquets faut-il ?
\end{exercice}\end{blocLivre}
\begin{blocLivre}{ex-7-13}\begin{exercice}[3]{Un pourcentage inverse}
Après une remise de 25\,\%, un objet coûte 54 euros. Retrouver son prix avant remise.
\end{exercice}\end{blocLivre}
\begin{blocLivre}{ex-7-14}\begin{exercice}[3]{Un arc de cercle}
Un cercle de rayon 6 cm porte un arc de mesure $75^\circ$. Calculer la longueur de cet arc, exactement puis au dixième de cm.
\label{fig-ex-7-14}\figureGeo{arc-75}
\end{exercice}\end{blocLivre}
\begin{blocLivre}{ex-7-15}\begin{exercice}[3]{Deux groupes réunis}
Un club compte 20 jeunes dont 40\,\% de filles, et 30 adultes dont 60\,\% de femmes. Quel est le pourcentage de personnes de sexe féminin dans le club ?
\end{exercice}\end{blocLivre}

\clearpage
